% Independently checked structural portion of the ordered-vector candidate.
% Amenability and the Ramsey expansion are not asserted in this note.
\begin{proposition}[Proved: the structural hypotheses for punctured ordered vector spaces]
\label{A:ordered-vector-structure}
Fix a finite field $k$. Let $\mathcal V_k$ consist of finite-dimensional
$k$-vector spaces with zero omitted from the universe, equipped with an
arbitrary independent linear order. Include the empty structure for the
zero vector space. In the relational language use one relation for each
linear equation on nonzero vectors; require the strong embeddings to
extend to injective linear maps. Then (S), (Rig), (SA), (FA),
(FA$\varnothing$), (St), (Op), (WEI), and (Fin) hold. Its category of
nonempty objects is total.
\end{proposition}
\begin{proof}
There are countably many finite types. Strong subobjects are exactly
punctured linear subspaces with the induced order, and compositions and
the intermediate-substructure axiom in (S) hold. Pushouts of vector
spaces provide amalgamation with the two factors meeting exactly in the
base. Orders on their union extend because the two orders agree on
the intersection, and the other vectors may be placed afterward.
The linear order makes every finite object rigid.

We give a functorial version of that construction. For fixed
$P\leq Z$, select a vector-space complement $Z_0$ with $Z=P\oplus Z_0$.
For every $X\geq P$, take $X\oplus Z_0$. The union of its embedded
$X$ and $Z$ is ordered by putting the $X$-side before the $Z$-side
in each gap of the order on $P\setminus\{0\}$. The remaining vectors
are the pairs $(x,z)$ with $x\in X\setminus P$ and $0\ne z\in Z_0$.
Put all of these after the union, ordered first by any fixed order on
$Z_0\setminus\{0\}$ and then by the given order on $X\setminus P$.
For an increasing linear embedding $u:X\to Y$ fixing $P$, the map
$u\oplus\mathrm{id}$ preserves all three parts and their orders:
$u^{-1}(P)=P$, so no mixed vector becomes an old vector. It is strong
and the construction respects composition. This proves (FA), including
the case $P=0$, which is (FA$\varnothing$) in the punctured language.

Here is the finite bridge needed for (St). Let $B,C\leq D$ be finite
subspaces, $F=B\cap C$, and let $d_0,d_1:D\to E$ be prescribed
increasing linear embeddings agreeing on $F$. Take an even number
$s\geq8$ of factors around a circle, one factor $E$ and $s-1$ copies
of $D$. Identify alternate boundary subspaces $B,C$, using $d_0,d_1$
at the two boundaries of $E$. Form the vector-space colimit $Z_0$:
the direct sum of all factors modulo the linear span of their overlap
relations. We check carefully that every factor embeds.

Index a boundary coefficient by an element of its overlap subspace.
A linear combination of overlap relations has, in each interior
$D$-factor, the difference of its two adjacent boundary coefficients.
If this combination is supported in a single factor $T$, then at every
other interior factor the coefficients agree as vectors of $D$, and
therefore belong to $B\cap C=F$. If $T=E$, all boundary coefficients
are the same element of $F$, and the two contributions in $E$ cancel
because $d_0,d_1$ agree on $F$. If $T$ is an interior factor, at least
one of the two chains between $T$ and $E$ has an interior factor, so
propagation along it produces a coefficient $f\in F$. If the other
chain has interior factors, it similarly produces $f'\in F$; the
vanishing coefficient at $E$ is $d_0(f)-d_1(f')$, so $f=f'$. If that
other chain has no interior factor, its coefficient $x$ initially
lies only in the appropriate boundary subspace, but the equation at
$E$ is $d_i(x)=d_{1-i}(f)=d_i(f)$, so again $x=f$. In either case
the two coefficients adjacent to $T$ coincide. The coefficient in
$T$ is then zero as well. Thus the
relation subspace meets every factor trivially, as required.

We also need the exact identifications between the old factor sets in
order to extend their linear orders. The same calculation with support
in two factors gives this description. Along either arc between those
factors, all intervening boundary coefficients are equal and lie in
$F$ as soon as the arc crosses an interior $D$-factor. Consequently two
vectors in different factors become equal precisely through the
prescribed boundary identifications, allowing the two endpoints of
$E$ to use the actual equality of their images in $E$. Equivalently,
the set union of the embedded factors is the strict set quotient used
in the alternating bridge: $F$-vectors have full visibility, and each
other vector has a proper visibility interval of length at most three.
For adjacent factors the long complementary arc contributes only an
$F$-vector and adds no identification; for factors at distance two the
short arc can only pass through $E$; at larger distance both arcs
contribute only $F$. This covers every pair of factors.

Embed the finite order on $E\setminus\{0\}$ into $\mathbb R$. On the
nonzero vectors of $D$ interpolate between the real coordinates of
$d_0$ and $d_1$. These increasing maps are constant on $F$. As in
Proposition~\ref{B:poset-bridge}, subdivide the interval finely and
alternate a map retaining the old $B$-coordinates with the next full
interpolated map; the intermediate map already has the new
$C$-coordinates. A sufficiently fine subdivision preserves every
strict order gap. The coordinates descend through exactly the set
identifications just proved, giving a linear order on the union of
the factor sets after arbitrary tie breaking. Put the remaining
nonzero vectors of $Z_0$ afterward in any order. All factors are now
increasing strong subspaces of a finite member of $\mathcal V_k$.
This is the required alternating bridge with prescribed endpoints.

If $M$ is the Fra\"isse limit and $g\in G_{B\cap C}$, apply that bridge
to the identity on a finite strong $D\supseteq B+C$ and to $g|_D$,
inside a finite strong $E\supseteq D+gD$. Embed the bridge over $E$
into $M$ and extend each copy of $D$ to an automorphism. Consecutive
automorphisms differ alternately by point stabilizers of $B$ and $C$.
Hence $g\in\langle G_B,G_C\rangle$, which proves (St).

Finite strong supports are closed under intersections. The
minimal-support argument therefore applies: for an open $U$, minimize
the dimension of $F$ with $G_F\leq U$. For every $u\in U$, (St) gives
$G_{F\cap uF}\leq U$, so minimality implies $uF=F$. Since $u$ preserves
the finite independent order on $F\setminus\{0\}$, it fixes that set
pointwise, and $U=G_F$. Thus (Op), and therefore (WEI), hold.

For clarity, (Fin) also follows without an implicit algebraic-closure
assumption. Strong amalgamation gives arbitrarily many independent
copies over a finite strong $F$ of a finite strong extension containing
any chosen point outside $F$. Thus every $G_F$-orbit outside $F$ is
infinite. If $G_C\leq G_F$ has finite index, each point of $C$ has a
finite $G_F$-orbit and must lie in $F$. Therefore $G_C\supseteq G_F$,
so they are equal. By (Op) this handles every pair of open subgroups.
Finally (FA), (FA$\varnothing$), and (Op) give totality by
Theorem~\ref{B:op-total}.
\end{proof}
